\section{systemd és szolgáltatáskezelés}

A \texttt{systemctl} és \texttt{journalctl} eszközök a szolgáltatások és naplók kezelésének alapjai systemd rendszereken.

\subsection{Alap műveletek}
\begin{lstlisting}[language=bash]
# Start
systemctl start httpd
# Stop
systemctl stop httpd
# Restart
systemctl restart httpd
# Enable at boot
systemctl enable httpd
# Check status
systemctl status httpd
\end{lstlisting}

\subsection{Naplók}
\begin{lstlisting}[language=bash]
# Logs for a service
journalctl -u httpd -n 200
# Follow in real time
journalctl -f
# Time interval
journalctl --since "2026-05-01" --until "2026-05-08"
\end{lstlisting}

Tippek:
\begin{itemize}
  \item Használjunk persistens journalt (ha szükséges) és forgassuk a helyi naplókat.
  \item A naplókban található érzékeny információt kezeljük GDPR/ellenőrzési szempontok szerint.
\end{itemize}
